\documentclass[12pt,a4paper]{article}
\usepackage[utf8]{inputenc}
\usepackage[spanish,german]{babel}
\usepackage[T1]{fontenc}
\usepackage{geometry}
\geometry{a4paper, margin=2.5cm}
\usepackage{titlesec}
\usepackage{enumitem}
\usepackage{multicol}
\usepackage{csquotes}

\titleformat{\section}
{\Large\bfseries}
{}
{0em}
{}[\titlerule]

\titleformat{\subsection}
{\large\bfseries}
{}
{0em}
{}

\setlength{\parindent}{0pt}
\setlength{\parskip}{0.5em}

\begin{document}

\title{\textbf{DIALOG: Saludo y Presentación}}
\author{Spanisch A1/A2}
\date{}
\maketitle

\textbf{Tema:} Saludo y presentación \\
\textbf{Nivel:} A1 \\
\textbf{Complejidad:} Einfach \\
\textbf{Palabras:} 687

\vspace{1em}

\textit{Nota: Se recomienda revisar primero el vocabulario antes de leer el diálogo.}

\vspace{1em}

\section*{Diálogo}

\textbf{María:} ¡Hola! ¿Cómo estás?

\textbf{Carlos:} ¡Hola! Muy bien, gracias. ¿Y tú? ¿Cómo estás?

\textbf{María:} Estoy muy bien, gracias. Me llamo María. ¿Cómo te llamas?

\textbf{Carlos:} Mucho gusto, María. Yo me llamo Carlos. ¿De dónde eres?

\textbf{María:} Soy de Buenos Aires, Argentina. ¿Y tú? ¿De dónde eres?

\textbf{Carlos:} Yo soy de Bogotá, Colombia. ¿Cuántos años tienes?

\textbf{María:} Tengo veinticinco años. ¿Y tú?

\textbf{Carlos:} Tengo veintisiete años. ¿Qué estudias o qué trabajas?

\textbf{María:} Estudio medicina en la universidad. Me gusta mucho ayudar a las personas. ¿Y tú? ¿Qué haces?

\textbf{Carlos:} Trabajo en una empresa de tecnología. Soy programador. Me encanta mi trabajo.

\textbf{María:} ¡Qué interesante! ¿Dónde vives ahora?

\textbf{Carlos:} Vivo en el centro de Bogotá, cerca del trabajo. ¿Y tú? ¿Dónde vives?

\textbf{María:} Vivo en un apartamento pequeño en Buenos Aires, cerca de la universidad. Es muy cómodo.

\textbf{Carlos:} ¿Tienes familia aquí?

\textbf{María:} Sí, mis padres viven en Buenos Aires también. Tengo una hermana menor. Se llama Ana. ¿Y tú? ¿Tienes familia?

\textbf{Carlos:} Sí, tengo dos hermanos. Uno es mayor y otro es menor que yo. Mis padres viven en Bogotá también.

\textbf{María:} ¿Qué te gusta hacer en tu tiempo libre?

\textbf{Carlos:} Me gusta leer libros, ver películas y salir con amigos. También me gusta hacer deporte. ¿Y a ti? ¿Qué te gusta hacer?

\textbf{María:} A mí me gusta mucho la música. Toco la guitarra. También me gusta cocinar y pasear por el parque.

\textbf{Carlos:} ¡Qué bien! Tal vez algún día podemos tocar música juntos. ¿Te gustaría?

\textbf{María:} ¡Sí, me encantaría! Eso sería muy divertido. ¿Tienes celular? Podemos quedar otro día.

\textbf{Carlos:} Sí, claro. Mi número es seis, cinco, cuatro, tres, dos, uno, cero, nueve, ocho, siete. ¿Cuál es tu número?

\textbf{María:} Mi número es nueve, uno, dos, tres, cuatro, cinco, seis, siete, ocho, cero. Llámame cuando quieras.

\textbf{Carlos:} Perfecto. Mucho gusto conocerte, María. Espero verte pronto.

\textbf{María:} Igualmente, Carlos. Fue un placer conocerte. ¡Hasta pronto!

\textbf{Carlos:} ¡Hasta luego! Que tengas un buen día.

\textbf{María:} Gracias, tú también. ¡Adiós!

\newpage

\section*{Vocabulario}

\begin{multicols}{2}
\begin{itemize}[leftmargin=*]
    \item \textbf{llamarse} - sich nennen, heißen
    \item \textbf{mucho gusto} - sehr erfreut, angenehm
    \item \textbf{¿De dónde eres?} - Woher kommst du?
    \item \textbf{estudiar} - studieren
    \item \textbf{medicina} - Medizin
    \item \textbf{universidad} - Universität
    \item \textbf{empresa} - Firma, Unternehmen
    \item \textbf{tecnología} - Technologie
    \item \textbf{programador} - Programmierer
    \item \textbf{encantar} - begeistern, lieben (me encanta = ich liebe)
    \item \textbf{apartamento} - Wohnung, Apartment
    \item \textbf{celular} - Handy, Mobiltelefon (südamerikanisch)
    \item \textbf{cómodo} - bequem, komfortabel
    \item \textbf{hermana menor} - jüngere Schwester
    \item \textbf{hermano mayor} - älterer Bruder
    \item \textbf{tiempo libre} - Freizeit
    \item \textbf{tocar} - spielen (Instrument), berühren
    \item \textbf{guitarra} - Gitarre
    \item \textbf{cocinar} - kochen
    \item \textbf{pasear} - spazieren gehen
    \item \textbf{quedar} - sich treffen, verabreden
    \item \textbf{igualmente} - ebenfalls, gleichfalls
    \item \textbf{placer} - Vergnügen, Freude
    \item \textbf{¡Hasta luego!} - Bis später!
    \item \textbf{¡Hasta pronto!} - Bis bald!
\end{itemize}
\end{multicols}

\newpage

\section*{Preguntas de Comprensión}

\begin{enumerate}[leftmargin=*]
    \item ¿Cómo se llama la chica en el diálogo?
    \item ¿De dónde es María?
    \item ¿Cuántos años tiene Carlos?
    \item ¿Qué estudia María?
    \item ¿Dónde trabaja Carlos?
    \item ¿Qué profesión tiene Carlos?
    \item ¿Dónde vive María?
    \item ¿Cuántos hermanos tiene Carlos?
    \item ¿Cómo se llama la hermana de María?
    \item ¿Qué le gusta hacer a María en su tiempo libre?
    \item ¿Qué instrumento toca María?
    \item ¿Qué le gusta hacer a Carlos en su tiempo libre?
    \item ¿Qué propone Carlos hacer con María?
    \item ¿Intercambian los dos sus números de teléfono?
    \item ¿Dónde vive Carlos?
    \item ¿Tiene María familia en Buenos Aires?
    \item ¿Qué edad tiene María?
    \item ¿Qué le gusta hacer a María además de tocar la guitarra?
    \item ¿Cómo se despiden al final del diálogo?
    \item ¿Qué quiere decir \enquote{mucho gusto}?
\end{enumerate}

\newpage

\section*{Verbo Irregular: SER (sein)}

\textit{Nota: \enquote{vosotros/as} se usa solo en España, no en español sudamericano. Se incluye aquí para completitud.}

\begin{multicols}{2}

\subsection*{Presente}
\begin{itemize}[leftmargin=*]
    \item yo - soy
    \item tú - eres
    \item él/ella/usted - es
    \item nosotros/as - somos
    \item vosotros/as - sois
    \item ustedes - son
    \item ellos/ellas - son
\end{itemize}

\subsection*{Pretérito Indefinido}
\begin{itemize}[leftmargin=*]
    \item yo - fui
    \item tú - fuiste
    \item él/ella/usted - fue
    \item nosotros/as - fuimos
    \item vosotros/as - fuisteis
    \item ustedes - fueron
    \item ellos/ellas - fueron
\end{itemize}

\subsection*{Pretérito Imperfecto}
\begin{itemize}[leftmargin=*]
    \item yo - era
    \item tú - eras
    \item él/ella/usted - era
    \item nosotros/as - éramos
    \item vosotros/as - erais
    \item ustedes - eran
    \item ellos/ellas - eran
\end{itemize}

\columnbreak

\subsection*{Futuro Simple}
\begin{itemize}[leftmargin=*]
    \item yo - seré
    \item tú - serás
    \item él/ella/usted - será
    \item nosotros/as - seremos
    \item vosotros/as - seréis
    \item ustedes - serán
    \item ellos/ellas - serán
\end{itemize}

\subsection*{Pretérito Perfecto}
\begin{itemize}[leftmargin=*]
    \item yo - he sido
    \item tú - has sido
    \item él/ella/usted - ha sido
    \item nosotros/as - hemos sido
    \item vosotros/as - habéis sido
    \item ustedes - han sido
    \item ellos/ellas - han sido
\end{itemize}

\end{multicols}

\newpage

\section*{Explicación Gramatical}

\textbf{Das Verb \enquote{gustar} -- Gefallen ausdrücken}

Im Spanischen funktioniert das Verb \enquote{gustar} (gefallen, mögen) anders als im Deutschen. Es wird mit indirekten Objektpronomen verwendet und das Subjekt steht im Singular oder Plural, je nachdem, was gefällt.

Die Struktur ist: \textbf{Indirektes Objektpronomen + gustar + Subjekt}

\textbf{Beispiele aus dem Text:}
\begin{itemize}[leftmargin=*]
    \item \enquote{Me gusta mucho ayudar a las personas} -- Mir gefällt es sehr, Menschen zu helfen
    \item \enquote{Me encanta mi trabajo} -- Ich liebe meine Arbeit (encantar funktioniert wie gustar)
    \item \enquote{¿Qué te gusta hacer en tu tiempo libre?} -- Was gefällt dir in deiner Freizeit zu tun?
    \item \enquote{A mí me gusta mucho la música} -- Mir gefällt Musik sehr (das \enquote{a mí} betont die Person)
\end{itemize}

Die indirekten Objektpronomen sind: me (mir), te (dir), le (ihm/ihr/Ihnen), nos (uns), les (ihnen/Ihnen). (Hinweis: In südamerikanischem Spanisch wird \enquote{os} nicht verwendet, stattdessen \enquote{les} für die 2. Person Plural.)

Wenn mehrere Dinge gefallen, verwendet man \enquote{gustan} (Plural): \enquote{Me gustan los libros} (Mir gefallen die Bücher).

\end{document}
